\section{Reference contract and execution model}
\label{sec:contract}
Let $\rho$ identify model and tokenizer bytes, prompt construction, numerical runtime, ordinary execution recipe, sampler and stopping rule. For input $x$, define the native endpoint
\[
\mathcal{O}_{\rho}(x)=(\mathbf y,\mathbf e,n,y_{\mathrm{last}},s),
\]
where $\mathbf y$ contains the full token-ID sequence, $\mathbf e$ its end-of-generation flags, $n$ the length and $s$ the stop reason. Qualification requires equality of every registered endpoint and all required evidence to be present. A shorter divergent response cannot qualify a latency comparison.

Three checks are kept separate. Exact-arithmetic equivalence concerns the mathematical operation. Bitwise operator correspondence compares outputs for fixed operands and execution conditions. Endpoint agreement compares the complete native result. Backend admission and actual placement are additional requirements: a mathematically valid view may be rejected by a backend, and an admitted view may still follow a different floating-point path.

A speculative transaction maintains separate target and draft state. The draft proposes up to $k$ tokens; the target evaluates the pending token and proposal positions; acceptance, correction or a bonus token determines the committed prefix. Rejected positions must be reconciled before the next cycle. A token that has been emitted but not yet consumed into KV remains pending and must be evaluated exactly once.

The query width $q$, caller width, microbatch capacity $U$ and physical graph shape are not interchangeable. A caller can contain multiple rows but execute through multiple physical graphs. We therefore report actual query extents, visible and physical key extents, grouped-head geometry, dtype, strides, masks and device placement when available. A local operator pass does not imply a full transaction pass; a full transaction pass on a finite panel does not establish general equivalence.
